\chapter{Proceso de desarrollo}
\section{Enfoque seleccionado}
Se propone un proceso iterativo, incremental y adaptativo. Cada Sprint de dos semanas debe tener un objetivo demostrable y un subconjunto de historias ajustado a la capacidad real. Scrum organiza la coordinación; Kanban hace visible el flujo; las prácticas técnicas de XP apoyan pruebas, refactorización e integración cuando existan las condiciones.

Esta es una adaptación para un trabajo académico, no una certificación de que ya se ejecutó Scrum. No se registran Daily Scrum, Sprint Review, retrospectivas, programación en pareja o TDD como actividades realizadas sin evidencia.

\section{Artefactos y eventos propuestos}
El Product Backlog mantiene el alcance candidato ordenado. El Sprint Backlog reúne el objetivo, historias seleccionadas y tareas. El incremento debe cumplir la Definition of Done acordada. Se propone una planificación inicial, un seguimiento breve y periódico, una revisión de producto y una retrospectiva al cierre; el equipo debe registrar fechas y resultados reales.

La guía detallada se conserva en \path{docs/PROCESO_DESARROLLO.md}. Los nombres de responsables propuestos no equivalen a una aceptación de sus funciones ni a una contribución ya realizada.

\section{Flujo y límites de trabajo}
Se plantea el tablero Pendiente, Lista, En curso, En revisión, Verificada y Hecha. Como política inicial para cuatro integrantes se recomienda una tarea de implementación activa por persona y un máximo de dos tareas en revisión. La coordinación puede reducir el WIP si aparecen bloqueos o archivos compartidos. Estos límites deben revisarse con observaciones del flujo, no tratarse como datos históricos.

Un bloqueo se anota con causa, responsable de resolverlo y siguiente acción. Si compromete el objetivo del Sprint se comunica al coordinador y se negocia alcance. No se declara terminada una tarea bloqueada por falta de acceso, revisión o integración.

\section{Definition of Done}
Una funcionalidad solo puede declararse terminada si cumple aceptación, valida entradas, persiste correctamente, aprueba las pruebas pertinentes, no tiene errores conocidos que impidan su flujo principal, cuenta con revisión humana registrada, está documentada y se integró en la rama acordada. La evidencia debe poder reproducirse.

Los criterios de una historia son específicos: por ejemplo, rechazar un correo duplicado. La DoD es transversal: por ejemplo, exigir revisión de código y documentación. Una batería local aprobada no satisface por sí sola todos los puntos de la DoD.

\section{Métricas y uso de IA}
Se proponen tiempo de ciclo, WIP, tareas bloqueadas, defectos encontrados y porcentaje de aceptación verificada. No se inventa velocidad de Sprints anteriores. La estimación de puntos es relativa y debe contrastarse con el equipo.

La plataforma identifica al agente como Codex basado en GPT-6; no expone aquí una variante exacta verificable. Se registra asistencia de IA sin atribuir autoría humana exclusiva. Las personas revisan requisitos, decisiones, código y evidencia. Los resultados de herramientas deben conservarse, y ningún texto generado sustituye una aprobación de Paz Sur.
